\textbf{结论名称：} 切线长定理

\textbf{适用条件：} 从圆外一点引圆的两条切线。

\textbf{变量说明：} 设 $\odot O$ 的半径为 $r$，点 $P$ 在圆外，$PA$、$PB$ 分别切圆于 $A$、$B$，则 $PA$、$PB$ 称为切线长。

\textbf{核心公式：} \[ \boxed{PA = PB,\quad \angle APO = \angle BPO} \]

\textbf{使用场景：} 证明线段相等；与勾股定理结合求切线长；利用角平分关系求解角度问题。

\textbf{简单例子：} 已知 $\odot O$ 半径 $r=3$，$OP=5$，$PA$ 为切线，则 $OA\perp PA$，\\
$PA=\sqrt{OP^{2}-OA^{2}}=\sqrt{5^{2}-3^{2}}=4$。\\
由定理得另一条切线长 $PB=4$；若两切线夹角 $\angle APB=60^{\circ}$，则 $\angle APO=\angle BPO=30^{\circ}$。